%!TeX root=main-DS1.tex
\def\Consist{\mathrm{Con}}
\chapter{Predikatski ra\v cun}

Iskazni ra\v cun nema dovoljnu izra\v zajnu mo\'c. Recimo, ako bismo slede\'ci iskaz
\begin{center}
  \textit{\enquote{Svaki prost broj koji je ve\'ci ili jednak sa 5}\\
  \textit{ima oblik $6k + 1$ ili $6k + 5$ za neko $k$}}
\end{center}
poku\v sali da zapi\v semo formulom iskaznog ra\v cuna, najbolje \v sto mo\v zemo da dobijemo je:
$$
  p \land q \Rightarrow r \lor s,
$$
gde smo slovima $p$, $q$, $r$ i $s$ ozna\v cili slede\'ce iskaze:

\begin{tabular}{l@{\quad$\leadsto$\quad}l}
  $p$ & \enquote{$n$ je prost broj}\\
  $q$ & \enquote{$n \ge 5$}\\
  $r$ & \enquote{$n = 6k + 1$}\\
  $s$ & \enquote{$n = 6k + 5$}.
\end{tabular}

\noindent
Lako se vidi da u formulaciji istog stava govornim jezikom ima mnogo vi\v se finesa, i da je zato
gornja iskazna formula veoma gruba reprezentacija tog iskaza.

Krajem XIX veka Gotlob Frege (Gottlob Frege) predla\v ze novi jezik matematike koji se danas zove \emph{predikatski ra\v cun}.
Frege oboga\'cuje iskazni ra\v cun tako \v sto uvodi na\v cin da se govori o svojstvima objekata
koriste\'ci posebne konstrukte koje zovemo \emph{predikati}. Predikati predstavljaju \enquote{iskazna slova sa strukturom},
na primer ovako:

\medskip

\begin{tabular}{l@{\quad$\leadsto$\quad}l}
  $P(n)$ & \enquote{$n$ je prost broj}\\
  $Q(n)$ & \enquote{$n \ge 5$}\\
  $R_1(n, k)$ & \enquote{$n = 6k + 1$}\\
  $R_5(n, k)$ & \enquote{$n = 6k + 5$}
\end{tabular}

\medskip

\noindent
i u jezik uvodi \emph{kvantifikatore:}

\medskip

\begin{tabular}{l@{\quad$\leadsto$\quad}l}
  $\forall$ & za svaki (od nema\v ckog \enquote{alle}), i\\
  $\exists$ & postoji (od nema\v ckog \enquote{existiert}).\\
\end{tabular}

\medskip

\noindent
U ovom novom jeziku mogu\'ce je iskazati mnoge matemati\v cke tvrdnje, pa i onu sa po\v cetka:
$$
  (\forall n) (P(n) \land Q(n) \Rightarrow (\exists k) (R_1(n, k) \lor R_5(n, k)).
$$
Sada je struktura govorne re\v cenice mnogo bolje predstavljena formulom koja se \v cita ovako:
$$
  \begin{array}{cccccccccccccccccccc}
    (\forall n)          & \Big( & P(n)               & \land    & Q(n)    & \Rightarrow \\
    \text{Za svako } n,  &   & n\text{ je prost}  & \text{i} & n \ge 5 & \text{implicira}\\[3mm]
    & & (\exists k)          & ( & R_1(n, k)            & \lor     & R_5(n, k) & ) & \Big). \\
    & & \text{postoji } k \text{ takvo da }  &   & n = 6k + 1  & \text{ili} & n = 6k + 5
  \end{array}
$$

\section{Formule predikatskog ra\v cuna}

Zbog ve\'ce izra\v zajne mo\'ci predikatski ra\v cun je komplikovaniji od iskaznog, pa \'ce nam
za precizno uvo\dj enje pojma formule predikatskog ra\v cuna biti potrebno malo vi\v se vremena.

Predikati predikatskog ra\v cuna opisuju osobine i uzajamne odnose objekata o kojima govorimo.
Objekte ozna\v cavamo \emph{promenljivim}, dok me\dj usobne odnose objekata ozna\v cavamo
\emph{predikatima} koji povezuju neke od tih promenljivih.
Predikat ima \emph{ime} koje naj\v ce\v s\'ce zapisujemo kao latini\v cno slovo,
i ima \emph{arnost}, \v sto je broj objekata o kojima predikat govori. Na primer:

\begin{center}
\begin{tabular}{lll}
  \hline
  $P(x)$&neka osobina objekta $x$&predikat arnosti 1 \\
        &npr.\ \enquote{$x$ je paran broj} & (\emph{unarni predikat})\\
  \hline
  $Q(x, y)$&neki odnos objekata $x$ i $y$&predikat arnosti 2\\
        &npr.\ \enquote{$x$ je deljiv sa $y$} & (\emph{binarni predikat})\\
  \hline
  $R(x, y, m)$&neki odnos objekata $x$, $y$ i $m$&predikat arnosti 3\\
        &npr.\ \enquote{$x \equiv_m y$} & (\emph{ternarni predikat})\\
  \hline
  $S(x, y, z, t)$&neki odnos objekata $x$, $y$, $z$ i $t$&predikat arnosti 4\\
        &npr.\ \enquote{$|x - y| > |z - t|$} & (\emph{kvaternarni predikat})\\
  \hline
\end{tabular}
\end{center}

\noindent
Pre nego \v sto krenemo da gradimo formule, moramo se dogovoriti o tome koja slova koristimo kao promenljive
i koja svojsta objekata \v zelimo da modelujemo. Zato je potrebno specificirati 
spisak promenljvih, kao i \emph{jezik} na kome govorimo, a to je spisak predikata i njihovih arnosti.
Na primer,

\begin{center}
  \begin{tabular}{ll}
    Promenljive: & $x$, $y$, $z$, $u$, $v$\\
    Jezik:   & $A(\cdot)$, $B(\cdot)$, $P(\cdot, \cdot)$, $W(\cdot, \cdot, \cdot, \cdot, \cdot)$.
  \end{tabular}
\end{center}

\noindent
Ovaj jezik ima dva unarna predikata, jedan binarni predikat i jedan predikat arnosti pet.
U ovom zapisu ta\v ckicama su ozna\v cena prazna mesta na koja mo\v zemo da stavimo neku od promenljivih
$x$, $y$, $z$, $u$, $v$. U jeziku \'cemo uvek imati simbol $=$ koji ozna\v cava jednakost
i njega nikada ne\'cemo navoditi u spisku predikata jer se podrazumeva.

\begin{DEF}
  Neka je dat \emph{kona\v can niz promenljivih} $x_1, x_2, \ldots, x_h$ i jezik $L$.
  \emph{Atomi\v cka formula} na jeziku $L$ je niz simbola oblika
  $$
    x_i = x_j,
  $$
  gde su $x_i$ i $x_j$ neke od promenljivih $x_1$, $x_2$, $x_3$, \ldots, $x_h$,
  ili niz simbola oblika
  $$
    R(x_{i_1}, x_{i_2}, \ldots, x_{i_n}),
  $$
  gde je $R$ predikat jezika $L$ arnosti $n$, a $x_{i_1}$, $x_{i_2}$, \ldots, $x_{i_n}$
  su neke od promenljivih $x_1$, $x_2$, $x_3$, \ldots, $x_h$ me\dj u kojima mo\v ze biti i istih.
\end{DEF}

\begin{NAP}\it
  Obratimo pa\v znju na to da je u atomi\v ckoj formuli nanizano onoliko promenljivih kolika je arnost predikata!
\end{NAP}

Atomi\v cke formule su najjednostavnije predikatske formule. Od njih, logi\v ckih veznika i kvantifikatora pravimo
slo\v zenije predikatske formule kako je opisano u definiciji koja sledi.

\begin{DEF}
  Neka je dat kona\v can niz promenljivih $x_1, x_2, \ldots, x_h$ i jezik $L$.
  \emph{Predikatska formula} na jeziku $L$ je niz simbola koji je formiran prema slede\'cim pravilima:
  \begin{enumerate}\itemsep -1mm
  \item
    atomi\v cke formule na jeziku $L$ su predikatske formule;
  \item
    ako su $\phi$ i $\psi$ predikatske formule, onda su to i slede\'ci nizovi simbola:
    $$
      \lnot \phi, \quad (\phi \land \psi), \quad (\phi \lor \psi), \quad (\phi \Rightarrow \psi), \quad (\phi \Leftrightarrow \psi);
    $$
  \item
    ako je $\phi$ predikatska formula i $x_i$ neka promenljiva, onda su i slede\'ci nizovi simbola predikatske formule:
    $$
      (\forall x_i)\phi, \quad (\exists x_i)\phi
    $$
  \item
    svaka predikatska formula se dobija kona\v cnom primenom pravila (1)--(3).
  \end{enumerate}
  (Zahtev (4) obezbe\dj uje da predikatske formule budu uvek kona\v cni nizovi simbola, kao i da nijedan drugi niz
  simbola nije predikatska formula.)
\end{DEF}

\paragraph{Primer.} Slede\'ci nizovi simbola su predikatske formule nad jezikom $P(\cdot)$, $Q(\cdot, \cdot)$, $R(\cdot, \cdot, \cdot)$
i sa promenljivim $x$, $y$, $z$:
\begin{itemize}
\item $R(z, y, z)$,
\item $(\forall x) Q(x, y)$,
\item $(\forall x)(\exists y)(Q(x, y) \land P(y))$,
\item $(\forall x)\big((\exists y) Q(x, y) \Rightarrow (\forall z) R(z, y, x)\big)$.
\end{itemize}
Primetimo: svaka promenljiva koja se javlja u poslednje dve formule je \enquote{napadnuta} kvantifikatorom,
dok sa prve dve formule to nije slu\v caj: u prvoj formuli nijedna promenljiva nije \enquote{napadnuta} kvantifikatorom,
a u drugoj formuli promenljiva $x$ je napadnuta kvantifikatorom, dok promenljiva $y$ nije.

\begin{NAP}\it
U ovom ud\v zbeniku razmatra\'cemo uglavnom predikatske formule kod kojih je svaka promenljiva
\enquote{napadnuta} kvantifikatorom. Takve formule se zovu \emph{re\v cenice}.
\end{NAP}

\begin{CONVENTION}
  Kao i kod iskaznih formula, prilikom zapisivanja predikatskih formula usvajamo slede\'ce konvencije:
  \begin{itemize}
  \item
    nikada ne pi\v semo spolja\v snje zagrade kod formula koje ne po\v cinju kvantifikatorom; i
  \item
    veznici $\land$ i $\lor$ imaju vi\v si prioritet nego veznici~$\Rightarrow$ i~$\Leftrightarrow$.
  \item
    umesto $\lnot x = y$ \v cesto \'cemo pisati $x \ne y$.
  \end{itemize}
\end{CONVENTION}


\section{Interpretacija predikatske formule}

Posmatrajmo slede\'cu formulu sa promenljivim $x$ i $y$ i na jeziku $P(\cdot, \cdot)$, $Q(\cdot)$:
$$
  (\forall x)(\exists y)(P(x, y) \land Q(y)).
$$
Da li je ona ta\v cna?

\bigskip

Pitanje, zapravo, nema smisla jer ne znamo \v sta zna\v ce predikati $P$ i $Q$. Da bismo mogli da govorimo o
tome da li je formula ta\v cna moramo prvo da znamo:
\begin{enumerate}
\item
  sa kojim objektima radimo, i
\item
  \v sta zna\v ce predikati koji se javljaju u jeziku formule.
\end{enumerate}
Kada su poznata ova dva podatka, dobijamo strukturu koja se zove \emph{interpretacija formule}.

\begin{EX}
  Neka su $P$ i $Q$ predikati. Interpretirajmo ih na skupu $\NN$ prirodnih brojeva ovako:

  \medskip

  \begin{tabular}{l@{\quad$\leadsto$\quad}l}
    $P(x, y)$ & \enquote{$x < y$},\\
    $Q(x)$ & \enquote{$x$ je paran broj}.
  \end{tabular}
  
  \medskip
  
  \noindent
  U ovoj interpretaciji formula
  $$
    (\forall x)(\exists y)(P(x, y) \land Q(y))
  $$
  dobija slede\'ce zna\v cenje:
  \begin{center}
    \textit{za svaki prirodan broj $x$ postoji prirodan broj $y$ takav da je $x < y$ i $y$ je paran,}
  \end{center}
  \v sto je ta\v cno tvr\dj enje.
\end{EX}

\begin{EX}
  Neka su $P$ i $Q$ predikati. Interpretirjamo ih sada na skupu $\NN$ ovako:

  \medskip

  \begin{tabular}{l@{\quad$\leadsto$\quad}l}
    $P(x, y)$ & \enquote{$x > y$},\\
    $Q(x)$ & \enquote{$x$ je paran broj}.
  \end{tabular}
  
  \medskip
  
  \noindent
  U ovoj interpretaciji formula
  $$
    (\forall x)(\exists y)(P(x, y) \land Q(y))
  $$
  dobija slede\'ce zna\v cenje:
  \begin{center}
    \textit{za svaki prirodan broj $x$ postoji prirodan broj $y$ takav da je $x > y$ i $y$ je paran,}
  \end{center}
  \v sto nije ta\v cno tvr\dj enje, jer za $x = 1$ ne mo\v zemo da na\dj emo paran prirodan broj koji je manji od njega.
\end{EX}

\begin{NAP}\it
  Dakle, ta\v cnost formule zavisi od interpretacije predikata. Kao \v sto smo upravo videli, mo\v ze se desiti da jedna ista formula
  bude ta\v cna u jednoj interpretaciji, a neta\v cna u drugoj.
\end{NAP}

\ifBIOINF\else
    \begin{DEF}
      Neka je dat jezik $L$ koga \v cine simboli $Q_1, Q_2, \ldots$ \v cije arnosti su, redom, $n_1, n_2, \ldots$.
      \emph{Interpretaciju} jezika $L$ \v cine skup $D$ koji se zove \emph{domen interpretacije}, i relacije
      $\rho_1, \rho_2, \ldots$ na skupu $D$ \v cije arnosti su $n_1, n_2, \ldots$. Tada je relacija $\rho_1$
      interpretacija simbola $Q_1$, relacija $\rho_2$ je interpretacija simbola $Q_2$ i tako dalje.
      Dakle, svako relacijsko slovo jezika $L$ se interpretira kao odgovaraju\'ca relacija na domenu $D$ koja
      ima istu arnost kao njeno slovo.
    \end{DEF}
\fi

\begin{EX}
  Dat je jezik koga \v cini jedan binarni simbol $R(\cdot, \cdot)$. Na\'ci interpretaciju u kojoj je slede\'ca
  formula ta\v cna:
  \begin{align*}
    (\exists x)(\exists y)(\exists z)(\exists u)
    (&x \ne y \land x \ne z \land x \ne u \land y \ne z \land y \ne u \land z \ne u \land\\
     &\land R(x, y) \land R(x, z) \land R(x, u) \land R(y, z) \land \lnot R(y, u) \land \lnot R(z, u)).
  \end{align*}
\end{EX}

\noindent
\begin{minipage}{3in}
\begin{proof}[Re\v senje]
  Ako pa\v zljivo pogledamo, ova formula ka\v ze da postoje \v cetiri razli\v cita elementa koji su
  u nekim me\dj usobnim odnosima. Zato \'cemo za domen interpretacije uzeti skup
  $D = \{a, b, c, d\}$, a binarni simbol $R$ \'cemo interpretirati binarnom relacijom
  $\rho$ na skupu $D$ koja je prikazana crte\v zom pored.
\end{proof}
\end{minipage}
\hfill
\begin{minipage}{2.5in}
  \begin{center}
    \input ft07-r12.pgf
  \end{center}
\end{minipage}

% \begin{EX}
%   Dat je jezik koga \v cini jedan unarni simbol $P(\cdot)$ i jedan binarni simbol $Q(\cdot, \cdot)$.
%   Na\'ci interpretaciju u kojoj je slede\'ca formula ta\v cna:
%   \begin{align*}
%     (\exists x)P(x) \; &\land \; (\exists y)\lnot P(y) \; \land \; (\forall x)\big(P(x) \Rightarrow\\
%      &(\exists y_1)(\exists y_2)(y_1 \ne y_2 \land \lnot P(y_1) \land \lnot P(y_2) \land Q(x, y_1) \land Q(x, y_2))\big).
%   \end{align*}
% \end{EX}

% \medskip

% \noindent
% \textit{Re\v senje.}
%   Unarni predikat opisuje svojstva pojedina\v cnih objekata. Prema tome, ima\'cemo dve vrste objekata, one koje
%   imaju svojstvo $P$ i one koje ga nemaju. Formula, dalje, ka\v ze da postoji objekat koji ima svojstvo $P$, kao i
%   objekat koji ga nema, i pri tome za svaki objekat koji ima svojstvo $P$ postoje dva razli\v cita objekta koji nemaju svojstvo
%   $P$ i koji su sa tim objektom u vezi~$Q$.

%   Neka je $D = \{a, b, c, d\}$ domen interpretacije. Unarni simbol $P$ \'cemo interpretirati
%   unarnom relacijom $\pi = \{a, b\}$ na skupu $D$, a binarni simbol $Q$ \'cemo interpretirati binarnom relacijom $\rho$
%   na skupu~$D$. 

% \smallskip

% \noindent
% \begin{minipage}{3in}
%   \hspace*{1.5em}Elementi binarne relacije $\rho$ su na crte\v zu pored ozna\v ceni strelicama, dok su elementi
%   unarne relacije $\pi$ ozna\v ceni crnim ta\v ckicama. U ovoj interpretaciji formula ka\v ze: \enquote{Postoji bar jedna crna
%   ta\v ckica i postoji bar jedna bela ta\v ckica, i za svaku crnu ta\v ckicu postoje dve razli\v cite bele ta\v ckice
%   koje su povezane strelicom sa tom crnom ta\v ckicom.}\hfill~$\square$
% \end{minipage}
% \hfill
% \begin{minipage}{2.5in}
%   \begin{center}
%     \input ft07-r13.pgf
%   \end{center}
% \end{minipage}

\begin{EX}
  Dat je jezik koga \v cini jedan ternarni simbol $R(\cdot, \cdot, \cdot)$.
  Na\'ci interpretaciju u kojoj je slede\'ca formula ta\v cna:
  $$
    (\forall x)(\forall y)(x \ne y \Rightarrow (\exists z)(z \ne x \land z \ne y \land R(x, z, y))).
  $$
\end{EX}

\noindent
\textit{Re\v senje.}
    Simbol $R$ \'cemo interpretirati relacijom \enquote{izme\dj u} na krugu kako sledi.
    Neka je $D = \{a, b, c\}$ domen interpretacije. Zamislimo da su ta\v cke $a$, $b$ i $c$ raspore\dj ene
    na krugu kako je to pokazano na slici ispod. Tada je svaka izme\dj u one druge dve i to samo treba formalno
    da zapi\v semo.

\medskip
    
\noindent
\begin{minipage}{3in}
    \hspace*{1.5em}U ovom slu\v caju \'ce biti najpogodnije da prosto navedemo sve ure\dj ene trojke
    koje \v cije ternarnu relaciju $\rho$ na skupu $D$ koja je interpretacija simbola~$R$:
    \begin{align*}
      \rho = \{ &(a, b, c), \, (c, b, a), \, (b, c, a),\\
                &(a, c, b), \, (c, a, b), \, (b, a, c)\}. \rlap{\hspace*{1cm}$\square$}
    \end{align*}
\end{minipage}
\hfill
\begin{minipage}{2.5in}
  \begin{center}
    \input ft07-r14.pgf
  \end{center}
\end{minipage}

\begin{DEF}
  Interpretacija u kojoj je fomula $\phi$ ta\v cna zove se \emph{model formule~$\phi$}.
  Interpretacija u kojoj je fomula $\phi$ neta\v cna zove se \emph{kontramodel formule~$\phi$}.
\end{DEF}

\begin{EX}
  Na\'ci model formule $(\forall x)(\forall y)(P(x, y) \Rightarrow (\exists z)(P(x, z) \land P(z, y)))$.
\end{EX}
\begin{proof}[Re\v senje]
  Interpretirajno na skupu $\RR$ realnih brojeva predikat $P$ ovako:
  \medskip

  \begin{tabular}{l@{\quad$\leadsto$\quad}l}
    $P(x, y)$ & \enquote{$x < y$}\\
  \end{tabular}
  
  \medskip
  
  \noindent
  U ovoj interpretaciji formula
  $$
    (\forall x)(\forall y)(P(x, y) \Rightarrow (\exists z)(P(x, z) \land P(z, y)))
  $$
  dobija slede\'ce zna\v cenje:
  \begin{center}
    \textit{za svaka dva realna broja $x$ i $y$, ako je $x < y$ onda}\\
    \textit{postoji realan broj $z$ takav da je $x < z$ i $z < y$}
  \end{center}
  \v sto je ta\v cno tvr\dj enje za realne brojeve (uvek mo\v zemo uzeti da je $z = \frac{x+y}{2}$).
\end{proof}

\begin{EX}
  Na\'ci kontramodel formule $(\forall x)(\forall y)(P(x, y) \Rightarrow (\exists z)(P(x, z) \land P(z, y)))$.
\end{EX}
\begin{proof}[Re\v senje]
  Da se podsetimo: kontramodel formule je interpretacija u kojoj formula nije ta\v cna.
  Interpretirajmo predikat $P$ na skupu $\NN$ ovako:

  \medskip

  \begin{tabular}{l@{\quad$\leadsto$\quad}l}
    $P(x, y)$ & \enquote{$x < y$}\\
  \end{tabular}
  
  \medskip
  
  \noindent
  U ovoj interpretaciji formula
  $$
    (\forall x)(\forall y)(P(x, y) \Rightarrow (\exists z)(P(x, z) \land P(z, y)))
  $$
  dobija slede\'ce zna\v cenje:
  \begin{center}
    \textit{za svaka dva \textbf{prirodna} broja $x$ i $y$, ako je $x < y$, onda}\\
    \textit{postoji \textbf{prirodan} broj $z$ takav da je $x < z$ i $z < y$,}
  \end{center}
  \v sto \emph{nije ta\v cno tvr\dj enje za prirodne brojeve} jer ne va\v zi za, recimo, $x = 1$ i $y = 2$.
\end{proof}

\begin{EX}
  $(a)$ Na\'ci predikatsku formulu koja ima slede\'cu osobinu: svaki model ove formule ima bar tri elementa.
  
  $(b)$ Neka je $n \ge 2$ prirodan broj.
  Na\'ci predikatsku formulu koja ima slede\'cu osobinu: svaki model ove formule ima bar $n$ elemenata.
\end{EX}
\begin{proof}[Re\v senje]
  $(a)$ To je formula: $(\exists x_1)(\exists x_2)(\exists x_3)(x_1 \ne x_2 \land x_1 \ne x_3 \land x_2 \ne x_3)$.
  
  $(b)$ To je formula: $(\exists x_1)(\exists x_2)\ldots (\exists x_n)(x_1 \ne x_2 \land x_1 \ne x_3 \land \ldots \land x_{n-1} \ne x_n)$.
\end{proof}

\begin{EX}
  $(a)$ Na\'ci predikatsku formulu koja ima slede\'cu osobinu: svaki model ove formule ima ta\v cno tri elementa.
  
  $(b)$ Neka je $n \ge 2$ prirodan broj.
  Na\'ci predikatsku formulu koja ima slede\'cu osobinu: svaki model ove formule ima ta\v cno $n$ elemenata.
\end{EX}
\begin{proof}[Re\v senje]
  $(a)$ To je formula:
  $$
    (\exists x_1)(\exists x_2)(\exists x_3)\big(x_1 \ne x_2 \land x_1 \ne x_3 \land x_2 \ne x_3 \land (\forall y)(y = x_1 \lor y = x_2 \lor y = x_3)\big).
  $$
  
  $(b)$ To je formula:
  \begin{align*}
    (\exists x_1)(\exists x_2)\ldots (\exists x_n)\big(&x_1 \ne x_2 \land x_1 \ne x_3 \land \ldots \land x_{n-1} \ne x_n \,\land\\
                                                       &(\forall y)(y = x_1 \lor y = x_2 \lor \ldots \lor y = x_n)\big). \qedhere
  \end{align*}
\end{proof}

Interesantno je da \emph{ne postoji} predikatska formula sa osobinom da svaki model te formule mora da bude kona\v can.
Tako\dj e, \emph{ne postoji} predikatska formula \v ciji svi modeli imaju paran broj elemenata.
Dokazi ova dva tvr\dj enja prevazilaze okvire ovog ud\v zbenika, ali nam daju va\v zan ose\'caj o ograni\v cenjima
izra\v zajne mo\'ci predikatskog ra\v cuna.


\section{Oblast dejstva kvantifikatora}

Zadatak kvantifikatora je da istakne neku promenljivu u formuli koja mu sledi
i da nam stavi do znanja da li bi formula treblo da va\v zi za svaku mogu\'cu vrednost te promenljive
(kvantifikator~$\forall$), ili za bar jednu njenu vrednost (kvantifikator~$\exists$).
Pri tome, kvantifikator deluje \emph{samo na formulu koja mu sledi!}
Na primer:
$$
  (\exists x)\underbrace{(P(x) \land Q(x))}_{\substack{\text{oblast dejstva}\\\text{kvantifikatora $\exists x$}}} \; \land \;
  (\exists y)\underbrace{(P(y) \land \lnot Q(y))}_{\substack{\text{oblast dejstva}\\\text{kvantifikatora $\exists y$}}}.
$$
Po\v sto dejstvo kvantifikatora prestaje kada se zavr\v si formula ispred koje se nalazi, \v cest je
obi\v caj da se isto slovo iskoristi vi\v se puta. Tako, gornja formula ima potpuno isto zna\v cenje kao i
slede\'ca formula:
$$
  \overset{K_1}{(\exists x)} \underbrace{(P(x) \land Q(x))}_{\substack{\text{oblast dejstva}\\\text{kvantifikatora $K_1$}}} \; \land \;
  \overset{K_2}{(\exists x)} \underbrace{(P(x) \land \lnot Q(x))}_{\substack{\text{oblast dejstva}\\\text{kvantifikatora $K_2$}}}.
$$

\begin{EX}
  Dat je jezik sa dva unarna simbola $P(\cdot)$ i $Q(\cdot)$. Posmatrajmo interpretaciju
  u kojoj na skupu $\NN$ prirodnih brojeva ovi predikati imaju slede\'ce zna\v cenje:

  \medskip

  \begin{tabular}{l@{\quad$\leadsto$\quad}l}
    $P(x)$ & \enquote{$x$ je paran broj},\\
    $Q(x)$ & \enquote{$x$ je neparan broj}.
  \end{tabular}
  
  \medskip

  \noindent
  Ta interpretacija je model formule
  $$
    (\exists \textcolor{red}{x})P(\textcolor{red}{x}) \land (\exists \textcolor{blue}{x})Q(\textcolor{blue}{x})
  $$
  jer postoji broj $\textcolor{red}{x} \in \NN$ koji je paran i postoji broj
  $\textcolor{blue}{x} \in \NN$ koji je neparan. Primetimo da, iako smo koristili isto slovo $x$
  i u levoj i u desnoj podformuli, kada se zavr\v si dejstvo prvog kvantifikatora promenljiva $x$ je
  \emph{oslobo\dj ena njegovog uticaja}, pa pod dejstvom drugog kvantifikatora mo\v ze da poprimi novu,
  mo\v zda neku druga\v ciju vrednost.

  Situacija je identi\v cna situaciji u programiranju gde je promenljiva vezana za blok u kome je deklarisana,
  pa kada se njen blok zavr\v si promenljiva je \emph{oslobo\dj ena} i mo\v ze da se koristi ponovo, sa drugim
  zna\v cenjem, u nekom drugom bloku. Na primer, u Javi mo\v ze da se napi\v se slede\'ce:
  \begin{verbatim}
    int[] domen = {1, 2, 3, 4, 5, 6, 7, 8, 9, 10};
    for(int x : domen){
        if(x % 2 == 0) System.out.println(x);
    }
    for(int x : domen){
        if(x % 2 == 1) System.out.println(x);
    }
  \end{verbatim}
\end{EX}


\begin{EX}
  Dat je jezik sa dva unarna simbola $P(\cdot)$ i $Q(\cdot)$. Na\'ci interpretaciju u kojoj slede\'ca
  formula nije ta\v cna:
  $$
    (\exists x)P(x) \land (\exists x)Q(x) \Rightarrow (\exists x)(P(x) \land Q(x)).
  $$
\end{EX}
\begin{proof}[Re\v senje]
  Neka su $P$ i $Q$ slede\'ci predikati na prirodnim brojevima:

  \medskip

  \begin{tabular}{l@{\quad$\leadsto$\quad}l}
    $P(x)$ & \enquote{$x$ je paran broj},\\
    $Q(x)$ & \enquote{$x$ je neparan broj}.
  \end{tabular}
  
  \medskip
  
  \noindent
  U ovoj interpretaciji formula
  $$
    (\exists x)P(x) \land (\exists x)Q(x) \Rightarrow (\exists x)(P(x) \land Q(x)).
  $$
  dobija slede\'ce zna\v cenje:
  \begin{center}
    \textit{ako postoji paran prirodan broj i postoji neparan prirodan broj,}\\
    \textit{onda postoji prirodan broj koji je i paran i neparan,}
  \end{center}
  \v sto nije ta\v cno tvr\dj enje.
\end{proof}

\begin{EX}
  Dat je jezik sa dva unarna simbola $P(\cdot)$ i $Q(\cdot)$. Na\'ci interpretaciju u kojoj slede\'ca
  formula nije ta\v cna:
  $$
    (\forall x)(P(x) \lor Q(x)) \Rightarrow (\forall x)P(x) \lor (\forall x)Q(x).
  $$
\end{EX}
\begin{proof}[Re\v senje]
  Neka su $P$ i $Q$ predikati interpretirani na prirodnim brojevima na slede\'ci na\v cin:

  \medskip

  \begin{tabular}{l@{\quad$\leadsto$\quad}l}
    $P(x)$ & \enquote{$x$ je paran broj},\\
    $Q(x)$ & \enquote{$x$ je neparan broj}.
  \end{tabular}
  
  \medskip
  
  \noindent
  U ovoj interpretaciji formula
  $$
    (\forall x)(P(x) \lor Q(x)) \Rightarrow (\forall x)P(x) \lor (\forall x)Q(x)
  $$
  dobija slede\'ce zna\v cenje:
  \begin{center}
    \textit{ako je svaki prirodan broj paran ili neparan,}\\
    \textit{onda su svi prirodni brojevi parni, ili su svi prirodni brojevi neparni,}
  \end{center}
  \v sto nije ta\v cno tvr\dj enje.
\end{proof}



\section{Valjane formule}

Predikatska formula mo\v ze u nekoj interpretaciji da bude ta\v cna, a u nekoj drugoj neta\v cna. Me\dj utim,
postoje predikatske formule koje su uvek ta\v cne, nezavisno od interpretacije.
Takve formule se zovu valjane formule.

\begin{DEF}
  \emph{Valjana formula} je re\v cenica koja je ta\v cna u svakoj interpretaciji.
  Izjavu \enquote{\textit{$\phi$ je valjana formula}} skra\'ceno zapisujemo ovako
  $$
    \vDash \phi.
  $$
\end{DEF}

\begin{EX}
  Od svake tautologije se mo\v ze napraviti valjana formula, na primer ovako:
  
  \medskip
  
  \begin{tabular}{ll}
    \hline
    Tautologija & Valjana formula\\
    \hline
    $p \lor \lnot p$ & $(\forall x)(P(x) \lor \lnot P(x))$\\
    $p \land (p \Rightarrow q) \Rightarrow q$ & $(\forall x)(P(x) \land (P(x) \Rightarrow Q(x)) \Rightarrow Q(x))$
  \end{tabular}
\end{EX}

Postoji mnogo valjanih formula koje ne poti\v cu od tautologija. U slede\'cem primeru
navodimo neke od najzna\v cajnijih.

\begin{EX}
  Dokazati da su slede\'ce formule valjane:

  \medskip
  
  \begin{tabular}{l}
  $\vDash (\forall x)(P(x) \land Q(x)) \Leftrightarrow (\forall x)P(x) \land (\forall x)Q(x)$\\[3mm]
  $\vDash (\exists x)(P(x) \lor Q(x))  \Leftrightarrow (\exists x)P(x) \lor  (\exists x)Q(x)$\\[3mm]
  $\vDash (\forall x)P(x) \lor (\forall x)Q(x) \Rightarrow (\forall x)(P(x) \lor Q(x))$\\[3mm]
  $\vDash (\exists x)(P(x) \land Q(x)) \Rightarrow (\exists x)P(x) \land (\exists x)Q(x)$\\[3mm]
  $\vDash \lnot(\exists x)P(x) \Leftrightarrow (\forall x)\lnot P(x)$\\[3mm]
  $\vDash \lnot(\forall x)P(x) \Leftrightarrow (\exists x)\lnot P(x)$
  \end{tabular}
\end{EX}
\begin{proof}[Re\v senje]
  Doka\v zimo da je $\phi = (\exists x)(P(x) \land Q(x)) \Rightarrow (\exists x)P(x) \land (\exists x)Q(x)$ valjana formula,
  a dokaz da su ostale valjane ostavljamo \v citaocu za ve\v zbu.
  
  Da bismo pokazali da je $\phi$ valjana formula
  treba da poka\v zemo da je ona ta\v cna u \emph{svakoj interpretaciji.}
  Neka je, zato, $D$ proizvoljan skup -- domen interpretacije. Unarni predikati $P$ i $Q$ se tada interpretiraju
  kao unarne relacije na skupu $D$. Unarna relacija na skupu $D$ je njegov proizvoljan podskup, pa \'cemo
  zato unarni predikat $P$ interpretirati kao unarnu relaciju $\rho \subseteq D$,
  a unarni predikat $Q$ kao unarnu relaciju $\sigma \subseteq D$.
  
  Formula $\phi$ je implikacija dve formule, pa \'cemo njenu ta\v cnost ustanoviti razmatranjem slede\'ca dva slu\v caja.
  
  Slu\v caj $1^\circ$: u interpretaciji $(D, \rho, \sigma)$ je formula $(\exists x)(P(x) \land Q(x))$ neta\v cna.
  Tada je u navedenoj interpretaciji cela formula $\phi$ ta\v cna (\emph{Ex falso quodlibet}) i dokaz je zavr\v sen.

  Slu\v caj $2^\circ$: u interpretaciji $(D, \rho, \sigma)$ je formula $(\exists x)(P(x) \land Q(x))$ ta\v cna.
  To zna\v ci da postoji $x \in D$ sa osobinom $P(x) \land Q(x)$, \v sto u interpretaciji $(D, \rho, \sigma)$
  zna\v ci da je $x \in \rho$ i $x \in \sigma$. Drugim re\v cima,
  postoji $x \in D$ takav da je $x \in \rho$, i postoji $x \in D$ takav da je $x \in \sigma$,
  \v sto je interpretacija formule $(\exists x)P(x) \land (\exists x)Q(x)$.
  Iako zvu\v ci kao obi\v cna jezi\v cka finesa, ovo je suptilan zaklju\v cak koji pokazuje da je
  formula $\phi$ ta\v cna u proizvoljnoj interpretaciji $(D, \rho, \sigma)$.
\end{proof}

\begin{EX}
  Dokazati da je slede\'ca formula valjana:
  $$
    (\exists x)(\exists y)(\forall z)(R(x, y, z) \lor \lnot R(y, x, z)).
  $$
\end{EX}
\begin{proof}[Re\v senje]
  Uo\v cimo proizvoljnu interpretaciju ove formule.
  Neka je $a$ proizvoljan objekt u toj interpretaciji i stavimo $x = y = a$. Tada formula postaje
  $$
    (\forall z)(R(a, a, z) \lor \lnot R(a, a, z)).
  $$
  Ova formula je ta\v cna nezavisno od toga \v sta zna\v ci predikat $R$ zato \v sto je iskaz
  $$
    R(a, a, z) \lor \lnot R(a, a, z)
  $$
  ta\v can za svaki objekat $z$ (Zakon isklju\v cenja tre\'ceg).
\end{proof}

Poka\v zimo sada kako se valjane formule koriste kao \enquote{transformaciona pravila}, kao \v sto je bio
slu\v caj sa tautologijama. Tako \'cemo dobiti alat koji \'ce nam omogu\'citi da transformi\v semo formule
u nove formule koje mo\v zda imaju jednostavniji ili \enquote{lep\v si} oblik, vode\'ci ra\v cuna o tome
da se pri transformacijama logi\v cka vrednost formule ne menja.

\ifBIOINF\else
    \begin{DEF}
      Predikatske formule $\phi$ i $\psi$ su \emph{ekvivalentne}, i pi\v semo $\phi \equiv \psi$, ako i samo ako je
      $\vDash \phi \Leftrightarrow \psi$.
    \end{DEF}
\fi 

Kao i ranije, simbol $\equiv$ ne treba me\v sati sa logi\v ckim veznikom $\Leftrightarrow$. Logi\v cki veznik
$\Leftrightarrow$ je deo jezika iskaznog ra\v cuna i koristi se kao simbol prilikom formiranja iskaznih formula, dok
simbol $\equiv$ pripada metajeziku.

\begin{EX}
  Dokazati da je $(\forall x)\phi \equiv \lnot (\exists x) \lnot \phi$.
\end{EX}
\begin{proof}[Re\v senje]
  Krenu\'cemo od formule sa desne strane i primeniti niz ekvivalentnih transformacija da bismo dobili formulu sa leve strane:
  \begin{align*}
    \lnot (\exists x) \lnot \phi
    &\;\;\equiv\;\; (\forall x) \lnot \lnot \phi\\
    &\;\;\equiv\;\; (\forall x) \phi
  \end{align*}
  jer je $\vDash \lnot \lnot \phi \Leftrightarrow \phi$ valjana formula izvedena iz odgovaraju\'ce tautologije.
\end{proof}

\begin{EX}
  Dokazati da je $(\exists x)\phi \equiv \lnot (\forall x) \lnot \phi$.
\end{EX}
\begin{proof}[Re\v senje]
  Krenu\'cemo od formule sa desne strane:
  \begin{align*}
    \lnot (\forall x) \lnot \phi
    &\;\;\equiv\;\; (\exists x) \lnot \lnot \phi\\
    &\;\;\equiv\;\; (\exists x) \phi\qedhere
  \end{align*}
\end{proof}

\v Cesto je potrebno na\'ci negaciju neke formule, i neophodne tehnike smo videli u slu\v caju
iskaznog ra\v cuna. U slu\v caju predikatskog ra\v cuna
\emph{negirati formulu $\phi$} zna\v ci krenuti od formule $\lnot \phi$ i primeniti niz ekvivalentnih
transformacija tako da na kraju dobijemo formulu u kojoj negacija napada samo atomi\v cke formule.

\begin{EX}\label{ft04.ex.density}
  Negirati formulu $(\forall x)(\forall y)(P(x, y) \Rightarrow (\exists z)(P(x, z) \land P(z, y)))$.
\end{EX}
\begin{proof}[Re\v senje]
  Ozna\v cimo datu formulu sa $\phi$. Tada:
  \begin{align*}
    \lnot \phi
    &\;\;\equiv\;\; \lnot(\forall x)(\forall y)(P(x, y) \Rightarrow (\exists z)(P(x, z) \land P(z, y)))\\
    &\;\;\equiv\;\; (\exists x)\lnot(\forall y)(P(x, y) \Rightarrow (\exists z)(P(x, z) \land P(z, y)))\\
    &\;\;\equiv\;\; (\exists x)(\exists y)\lnot(P(x, y) \Rightarrow (\exists z)(P(x, z) \land P(z, y)))\\
    &\;\;\equiv\;\; (\exists x)(\exists y)(P(x, y) \land \lnot(\exists z)(P(x, z) \land P(z, y)))\\
    &\;\;\equiv\;\; (\exists x)(\exists y)(P(x, y) \land (\forall z)\lnot(P(x, z) \land P(z, y)))\\
    &\;\;\equiv\;\; (\exists x)(\exists y)(P(x, y) \land (\forall z)(\lnot P(x, z) \lor \lnot P(z, y))).\qedhere
  \end{align*}
\end{proof}

Ovo je korisna ve\v stina izme\dj u ostalog i zato \v sto se nala\v zenje kontramodela formule
svodi na nala\v zenje modela njene negacije. Da se podsetimo,
kontramodel formule $\phi$ je interpretacija u kojoj formula $\phi$ nije ta\v cna.
Dakle, u toj interpretaciji mora biti ta\v cna formula $\lnot\phi$, pa se problem svodi na to da na\dj emo
model formule $\lnot\phi$:

\begin{NAP}\it
  Da bi se na\v sao kontramodel formule $\phi$ treba na\'ci model formule $\lnot \phi$.
\end{NAP}

\begin{EX}
  Na\'ci kontramodel formule
  $$
    \phi = (\forall x)(\forall y)((\exists z)(P(z, x) \land P(z, y)) \Rightarrow (\forall z)(P(z, x) \land P(z, y))).
  $$
\end{EX}
\begin{proof}[Re\v senje]
  Prema gornjoj napomeni, problem se svodi na to da na\dj emo model formule
  $$
    \lnot \phi = (\exists x)(\exists y)((\exists z)(P(z, x) \land P(z, y)) \land (\exists z)(\lnot P(z, x) \lor \lnot P(z, y))).
  $$
  Interpretirajmo predikat $P$ na skupu $\NN$ ovako:
  
  \medskip

  \begin{tabular}{l@{\quad$\leadsto$\quad}l}
    $P(x, y)$ & \enquote{$x < y$}.
  \end{tabular}
  
  \medskip
  
  \noindent
  Formula $\lnot\phi$ je ta\v cna, recimo, za $x = 2$ i $y = 3$ \v sto se lako vidi, jer je
  $1 < 2 \land 1 < 3$ i pri tome nije $4 < 3$.
\end{proof}

\ifBIOINF
  \endinput
\fi

\section{Uvo\dj enje novih kvantifikatora}

\v Ceste su situacije kada \v zelimo da ka\v zemo da \emph{postoji ta\v cno jedan objekat} sa nekom
osobinom~$\phi$. Formula
$$
  (\exists x) \phi(x)
$$
ka\v ze da objekat sa tom osobinom postoji, ali ne garantuje da je on jedinstven.
Kvantifikator $\exists$ mo\v ze da se unapredi do kvantifikatora $\exists!$ koga \v citamo
\enquote{postoji samo jedno}. Interesantno je da se ovim ne unosi nova izra\v zajna mo\'c u predikatski ra\v cun,
jer se formula
$$
  (\exists! x) \phi(x)
$$
(\enquote{postoji samo jedno $x$ sa osobinom $\phi(x)$}) mo\v ze shvatiti kao \enquote{skra\'cenica} za:
$$
  \underbrace{(\exists x) \phi(x)}_{\text{egzistencija}} \quad \land \quad
  \underbrace{(\forall y)(\forall z)(\phi(y) \land \phi(z) \Rightarrow y = z)}_{\text{jedinstvenost}}.
$$


Kada se govori o rezultatima koji se odnose na skupove, veoma su \v ceste izjave oblika
  \textit{Za svako $x \in A$ va\v zi \ldots}
ili
  \textit{Postoji $x \in A$ takvo da \ldots}
Zato je uobi\v cajeno da se jezik predikatskog ra\v cuna pro\v siri \emph{ograni\v cenim verzijama kvantifikatora $\forall$
i $\exists$} koje pi\v semo ovako:
\begin{center}
  \begin{tabular}{l@{\quad$\leadsto$\quad}l}
    $(\forall x \in A) \phi(x)$ & Za svako $x \in A$ va\v zi $\phi(x)$\\
    $(\exists x \in A) \phi(x)$ & Postoji $x \in A$ takvo da va\v zi $\phi(x).$
  \end{tabular}
\end{center}
Ni ograni\v ceni kvantifikatori ne unose novu izra\v zajnu mo\'c u predikatski ra\v cun ve\'c samo
predstavljaju pogodne \enquote{skra\'cenice}:
\begin{center}
  \begin{tabular}{l@{\quad je skra\'ceni zapis formule\quad}l}
    $(\forall x \in A) \phi(x)$ & $(\forall x)(x \in A \Rightarrow \phi(x))$ \\
    $(\exists x \in A) \phi(x)$ & $(\exists x)(x \in A  \land      \phi(x))$.
  \end{tabular}
\end{center}
Da bismo videli da je to tako i na formalnom nivou, pokaza\'cemo da se negiranje ograni\v cenih kvantifikatora pona\v sa
na o\v cekivani na\v cin.

\begin{THM}
  $(a)$ $\lnot (\exists x \in A) \phi(x) \equiv (\forall x \in A) \lnot \phi(x)$;

  $(b)$ $\lnot (\forall x \in A) \phi(x) \equiv (\exists x \in A) \lnot \phi(x)$.
\end{THM}
\begin{proof}[Dokaz]
  $(a)$ Koriste\'ci tautologiju $\models (p \Rightarrow q) \Leftrightarrow (\lnot p \lor q)$
  lako se vidi da je:
  \begin{align*}
    \lnot (\exists x \in A) \phi(x)
    & \equiv \lnot (\exists x)(x \in A \land \phi(x))\\
    & \equiv (\forall x)\lnot(x \in A \land \phi(x))\\
    & \equiv (\forall x)(\lnot(x \in A) \lor \lnot \phi(x))\\
    & \equiv (\forall x)(x \in A \Rightarrow \lnot \phi(x))\\
    & \equiv (\forall x \in A) \lnot \phi(x).
  \end{align*}
  Dokaz tvr\dj enja $(b)$ ostavljamo za ve\v zbu.
\end{proof}


% \section{Predikatski ra\v cun kao formalna teorija}

% U prethodnim poglavljima smo koristili \enquote{skra\'cenu definiciju} predikatskog ra\v cuna kako bismo
% pa\v znju \v citaoca fokusirali na odnos sintakse (formule) i semantike (interpretacije).
% Za nastavak diskusije nam treba puna snaga predikatskog ra\v cuna, pa \'cemo sada pro\v siriti
% definicije sa po\v cetka.

% \begin{DEF}
  % Jezik $\calL$ predikatskog ra\v cuna je unija tri po parovima disjunktna skupa:
  % $$
    % \calL = \calR \cup \calF \cup \calC,
  % $$
  % gde je $\calR$ skup \emph{relacionih simbola}, $\calF$ skup \emph{funkcijskih simbola} i $\calC$ skup
  % \emph{konstantnih simbola}. Svakom relacijskom i funkcijskom simbolu pridru\v zujemo prirodan broj koji
  % se zove \emph{arnost} tog simbola.
% \end{DEF}

% \begin{EX}
  % Na primer, jezik
  % $\calL = \{P(\cdot), Q(\cdot, \cdot, \cdot)\} \cup \{f(\cdot), g(\cdot, \cdot)\} \cup \{a, b\}$
  % ima dva relacijska simbola (simbol $P$ arnosti~1 i simbol $Q$ arnosti~3), dva funkcijska simbola
  % (simbol $f$ arnosti~1 i simbol $g$ arnosti~2) i dve konstante ($a$ i $b$).
% \end{EX}

% \begin{DEF}
  % Neka je dat kona\v can niz promenljivih $x_1, x_2, \ldots, x_h$ i jezik $\calL = \calR \cup \calF \cup \calC$.
  % \emph{Term} na jeziku $\calL$ je niz simbola izgra\dj en prema slede\'cim pravilima:
  % \begin{enumerate}\itemsep -1mm
  % \item
    % promenljive $x_1, x_2, \ldots, x_h$ i konstante iz $\calC$ su termi;
  % \item 
    % ako je $f \in \calF$ funkcijski simbol arnosti $n$ i $t_1, \ldots, t_n$ termi, onda je i niz simbola
    % $f(t_1, \ldots, t_n)$ term;
  % \item
    % svaka term se dobija kona\v cnom primenom pravila (1)--(2).
  % \end{enumerate}
  % (Zahtev (3) obezbe\dj uje da termi budu uvek kona\v cni nizovi simbola, kao i da nijedan drugi niz
  % simbola nije term.)
% \end{DEF}

% \begin{DEF}
  % Neka je dat kona\v can niz promenljivih $x_1, x_2, \ldots, x_h$ i jezik $\calL = \calR \cup \calF \cup \calC$.
  % \emph{Atomi\v cka formula} na jeziku $L$ je niz simbola oblika
  % $$
    % t_1 = t_2
  % $$
  % gde su $t_1$ i $t_2$ termi na jeziku $\calL$, ili niz simbola oblika
  % $$
    % R(t_1, t_2, \ldots, t_n)
  % $$
  % gde je $R \in \calR$ relacijski simbol jezika $\calL$ arnosti $n$, a $t_1$, $t_2$, \ldots, $t_n$
  % su termi na jeziku $\calL$.
% \end{DEF}

% \begin{DEF}
  % Neka je dat kona\v can niz promenljivih $x_1, x_2, \ldots, x_h$ i jezik $\calL = \calR \cup \calF \cup \calC$.
  % \emph{Predikatska formula} na jeziku $\calL$ je niz simbola koji je formiran prema slede\'cim pravilima:
  % \begin{enumerate}\itemsep -1mm
  % \item
    % atomi\v cke formule na jeziku $\calL$ su predikatske formule;
  % \item
    % ako su $\phi$ i $\psi$ predikatske formule, onda su to i slede\'ci nizovi simbola:
    % $$
      % \lnot \phi, \quad (\phi \land \psi), \quad (\phi \lor \psi), \quad (\phi \Rightarrow \psi), \quad (\phi \Leftrightarrow \psi);
    % $$
  % \item
    % ako je $\phi$ predikatska formula i $x_i$ neka promenljiva, onda su i slede\'ci nizovi simbola predikatske formule:
    % $$
      % (\forall x_i)\phi, \quad (\exists x_i)\phi
    % $$
  % \item
    % svaka predikatska formula se dobija kona\v cnom primenom pravila (1)--(3).
  % \end{enumerate}
  % (Zahtev (4) obezbe\dj uje da predikatske formule budu uvek kona\v cni nizovi simbola, kao i da nijedan drugi niz
  % simbola nije predikatska formula.)
% \end{DEF}

% U verziji koju \'cemo ovde pokazati predikatski ra\v cun kao formalna teorija ima \v sest \emph{shema-aksioma}
% i dva pravila izvo\dj enja. Shema-aksioma je modla po kojoj se dalje \v stancuju aksiome tako \v sto umesto \enquote{plejsholdera}
% $\phi$ i $\psi$ udenemo bilo koju formulu. To je na\v cin da se iz skupa pravila izbaci pravilo supstitucije
% (koje je ovde vrlo pipavo za uvo\dj enje), po cenu da aksiomatski sistem sada postaje \emph{beskona\v can skup
% formula} jer od svake shema-aksiome mo\v ze da se napravi beskona\v cno formula koje su zapravo aksiome sistema.

% Postojanje kvantifikatora i \v cinjenica da oni vezuju promenljive stvara \v citav niz tehni\v ckih pote\v sko\'ca.
% Izme\dj u ostalog, to je jedan od klju\v cnih razloga za\v sto smo odustali od pravila supstitucije i pre\v sli
% na shema-aksiome. Druga tehnikalija sa kojom \'cemo morati da se izborimo je re\v sena slede\'com definicijom.

% \begin{DEF}
  % Za term $t$ ka\v zemo da je \emph{slobodan za promenljivu $x$ u formuli $\phi(x)$} u kojoj se $x$ javlja kao slobodna
  % promenljiva, ako nakon zamene simbola $x$ termom $t$ nijedna promenljiva terma $t$ ne postaje vezana u formuli
  % $\phi(t)$.
% \end{DEF}

% \begin{EX}
  % Neka je $\phi(x) = (\forall y)R(x, y)$. Tada je term $t_1 = f(z, x)$ slobodan za $x$ u formuli $\phi(x)$
  % jer nakon zamene $x$ sa $t_1$ dobijamo
  % $\phi(t_1) = (\forall y)R(f(z, x), y)$ gde nijedna od promenljivih $z$ i $x$ nije napadnuta nijednim
  % kvantifikatorom formule $\phi$. S druge strane term $t_2 = f(y, x)$ nije slobodan za $x$ u formuli $\phi(x)$
  % jer nakon zamene $x$ sa $t_2$ dobijamo
  % $\phi(t_2) = (\forall y)R(f(y, x), y)$ gde je promenljiva $y$ u termu $f(y, x)$ napadnuta
  % kvantifikatorom formule $\phi$.
% \end{EX}

% Ako nam je fiksiran jezik $\calL$, predikatski ra\v cun kao formalna teorija je zadat slede\'cim shema-aksiomama:

% \begin{description}
% \item[A1:] $\phi \Rightarrow (\psi \Rightarrow \phi)$;
% \item[A2:] $(\phi \Rightarrow (\psi \Rightarrow \rho)) \Rightarrow ((\phi \Rightarrow \psi) \Rightarrow (\phi \Rightarrow \rho))$;
% \item[A3:] $(\lnot \psi \Rightarrow \lnot \phi) \Rightarrow (\phi \Rightarrow \psi)$;
% \item[A4:] $\lnot \phi \Rightarrow (\phi \Rightarrow \psi)$;
% \item[A5:] $(\forall x)(\phi \Rightarrow \psi(x)) \Rightarrow (\phi \Rightarrow (\forall x)\psi(x))$, ako promenljiva $x$
           % nije slobodna u formuli $\phi$;
% \item[A6:] $((\forall x)\phi(x)) \Rightarrow \phi(t)$ ako je term $t$ slobodan za promenljivu $x$ u  formuli $\phi$.
% \end{description}
% Pravila izvo\dj enja ove teorije su \emph{Modus ponens} i \emph{Pravilo generalizacije}:
% $$
  % \mathit{MP}: \begin{array}{c}
    % \phi \qquad \phi \Rightarrow \psi \\ \hline
    % \psi
  % \end{array},\qquad
  % \mathit{Gen}: \begin{array}{c}
    % \phi \\ \hline
    % (\forall x)\phi
  % \end{array}
% $$

% \medskip

% I u ovom pro\v sirenom ra\v cunu imamo pojam sintaksne posledice kao i
% izuzetno va\v znu i korisnu Teoremu dedukcije.

% \begin{DEF}
  % Neka je $\calL$ jezik predikatskog ra\v cuna,
  % neka je $\Sigma$ skup formula i neka je $\phi$ neka formula. Tada oznaka
  % $\Sigma \vdash \phi$ zna\v ci da postoji kona\v can niz formula
  % $$
    % \alpha_1, \; \alpha_2, \; \ldots, \; \alpha_n
  % $$
  % takav da je $\alpha_n = \phi$, i za svako $i \in \{1, \ldots, n\}$ formula $\alpha_i$ je ili aksioma
  % (dobijena od neke shema-aksiome zamenom \enquote{plejsholdera} proizvoljnim formulama), ili
  % neka od formula iz $\Sigma$, ili je dobijena primenom nekog pravila izvo\dj enja na neke od formula koje joj
  % prethode u nizu. Navedeni niz formula tada nazivamo \emph{dokaz formule $\phi$ iz pretpostavki $\Sigma$}.
% \end{DEF}

% \begin{THM}[Teorema dedukcije za predikatski ra\v cun]
  % Neka je $\Sigma$ skup formula predikatskog ra\v cuna na jeziku $\calL$
  % i neka su $\phi$ i $\psi$ jo\v s dve formule. Tada $\Sigma \vdash \phi \Rightarrow \psi$
  % ako i samo ako $\Sigma \union \{\phi\} \vdash \psi$, pri \v cemu nijedna od promenljivih pomo\'cu kojih
  % primenjujemo pravilo (Gen) nije slobodna u formuli $\phi$.
% \end{THM}

% Teorema dedukcije za predikatski ra\v cun dobija pitomiju formulaciju ako svedemo diskusiju na specijalnu klasu formula.

% \begin{DEF}
  % \emph{Re\v cenica} je predikatska formula bez slobodnih promenljivih.
% \end{DEF}

% \begin{THM}[Teorema dedukcije za predikatski ra\v cun V2.0]
  % Neka je $\Sigma$ skup formula predikatskog ra\v cuna na jeziku $\calL$
  % i neka su $\phi$ i $\psi$ jo\v s dve formule, pri \v cemu je $\phi$ re\v cenica.
  % Tada $\Sigma \vdash \phi \Rightarrow \psi$
  % ako i samo ako $\Sigma \union \{\phi\} \vdash \psi$.
% \end{THM}

% Sli\v cno iskaznom ra\v cunu, uz mnogo truda mo\v ze se pokazati da su sve valjane formule
% posledice formalne teorije predikatskog ra\v cuna, i nijedna druga formula to nije.
% Dakle, valjane formule mo\v zemo precizno opisati formalnim teorijama.

% \begin{THM}[G\"odel 1929]
  % Neka je $\phi$ re\v cenica. Tada $\vdash \phi$ ako i samo ako je $\phi$ valjana formula.
% \end{THM}

% Ako izjavu $\vDash \phi$ shvatimo kao \enquote{$\phi$ je ta\v cno}, a izjavu $\vdash \phi$ kao
% \enquote{$\phi$ je dokazivo iz aksioma} onda prethodna veoma va\v zna teorema mo\v ze da se shvati i ovako:
% formula predikatskog ra\v cuna je ta\v cna ako i samo ako je dokaziva iz aksioma.

% Dakle, i predikatski ra\v cun kako smo ga ovde pokazali ima tu super-va\v znu osobinu
% da se \emph{semantika} (pojam ta\v cnosti) i \emph{sintaksa} (pojam formalne dokazivosti) poklapaju,
% \v sto nije uvek slu\v caj!

% \section{Peanova aritmetika}

% Shema-aksiome A1--A6 se jo\v s zovu i \emph{logi\v cke aksiome} jer one reguli\v su
% deduktivni proces. Ako unutar predikatskog ra\v cuna \v zelimo da govorimo o nekoj specifi\v cnoj
% teoriji, onda je formalnom sistemu potrebno dodati specifi\v cne, \emph{ne-logi\v cke} aksiome
% koje reguli\v su specifi\v cno pona\v sanje.

% Na primer, Peanova aritmetika nam omogu\'cuje da rezonujemo o aritmetici kao formalnoj teoriji.
% Jezik Peanove aritmetike nema relacijskih simbola, ima jedan funkcijski simbol $s$ i jednu konstantu~0.
% Aksiome Peanove aritmetike se sastoje od dve \enquote{obi\v ne} aksiome i jedne shema-aksiome koje dodajemo
% predikatskom ra\v cunu kao \emph{ne-logi\v cke aksiome}:

% \begin{description}
% \item[PA1:] $\lnot(\exists x) s(x) = 0$;
% \item[PA2:] $(\forall x)(\forall y)(s(x) = s(y) \Rightarrow x = y)$;
% \item[PA3:] $\big(\phi(0) \land (\forall x)(\phi(x) \Rightarrow \phi(s(x)))\big) \Rightarrow (\forall x)\phi(x)$.
% \end{description}

% Skup aksioma Peanove aritmetike \'cemo ozna\v citi sa $PA$.

% \begin{EX}
  % Dokazati da $PA \vdash (\forall x)(x = 0 \lor (\exists y)x = s(y))$.

  % \medskip

  % \textit{Re\v senje.}
  % Neka je $\phi(x)$ formula $[x = 0 \lor (\exists y)x = s(y)]$. Formulu \'cemo ograditi uglastim zagradama
  % da bismo lak\v se mogli da ispratimo \v sta se de\v sava u dokazu.

% \begin{enumerate}
% \item%1
  % $\big([0 = 0 \lor (\exists y)0 = s(y)] \land (\forall x)([x = 0 \lor (\exists y)x = s(y)]
  % \Rightarrow [s(x) = 0 \lor (\exists y)s(x) = s(y)])\big) \Rightarrow (\forall x)[x = 0 \lor (\exists y)x = s(y)]$\newline
  % [PA3 raspisana za formulu $\phi$]
% \item%2
  % $0 = 0$\newline
  % [refleksivnost jednakosti]
% \item%3
  % $0 = 0 \Rightarrow [0 = 0 \lor (\exists y)0 = s(y)]$\newline
  % [tautologija $p \Rightarrow (p \lor q)$]
% \item%4
  % $[0 = 0 \lor (\exists y)0 = s(y)]$\newline
  % [MP na (2) i (3)]
% \item%5
  % $s(x) = s(x)$\newline
  % [refleksivnost jednakosti]
% \item%6
  % $s(x) = s(x) \Rightarrow (\exists y)s(x) = s(y)$\newline
  % [valjana formula $\phi(x, x) \Rightarrow (\exists y)\phi(x, y)$]
% \item%7
  % $(\exists y)s(x) = s(y)$\newline
  % [MP na (5) i (6)]
% \item%8
  % $((\exists y)s(x) = s(y)) \Rightarrow [s(x) = 0 \lor (\exists y)s(x) = s(y)]$\newline
  % [tautologija $p \Rightarrow (q \lor p)$]
% \item%9
  % $[s(x) = 0 \lor (\exists y)s(x) = s(y)]$\newline
  % [MP na (7) i (8)]
% \item%10
  % $[s(x) = 0 \lor (\exists y)s(x) = s(y)] \Rightarrow ([x = 0 \lor (\exists y)x = s(y)] \Rightarrow [s(x) = 0 \lor (\exists y)s(x) = s(y)])$\newline
  % [tautologija $p \Rightarrow (q \Rightarrow p)$]
% \item%11
  % $[x = 0 \lor (\exists y)x = s(y)] \Rightarrow [s(x) = 0 \lor (\exists y)s(x) = s(y)]$\newline
  % [MP na (9) i (10)]
% \item%12
  % $(\forall x)([x = 0 \lor (\exists y)x = s(y)] \Rightarrow [s(x) = 0 \lor (\exists y)s(x) = s(y)])$\newline
  % [Gen na (11)]
% \item%13
  % $(4) \Rightarrow ((12) \Rightarrow ((4) \land (12)))$\newline
  % [tautologija $p \Rightarrow (q \Rightarrow (p \land q))$; (4) i (12) su formule iz dokaznog niza]
% \item%14
  % $(12) \Rightarrow ((4) \land (12))$\newline
  % [MP na (4) i (13)]
% \item%15
  % $[0 = 0 \lor (\exists y)0 = s(y)] \land (\forall x)([x = 0 \lor (\exists y)x = s(y)] \Rightarrow [s(x) = 0 \lor (\exists y)s(x) = s(y)])$\newline
  % [MP na (12) i (14); ovo je formula $(4) \land (12)$]
% \item%16
  % $(\forall x)[x = 0 \lor (\exists y)x = s(y)]$\newline
  % [MP na (15) i (1)]
% \end{enumerate}
% Ovim je formalni dokaz zavr\v sen.
% \end{EX}

% Peanovoj aritmetici sada mo\v zemo lako da dodamo operacije sabiranja i mno\v zenja slede\'cim \emph{rekurzivnim definicijama}:

% \begin{description}
% \item[S1:] $x + 0 = x$;
% \item[S2:] $x + s(y) = s(x + y)$;
% \item[M1:] $x \cdot 0 = 0$;
% \item[M2:] $x \cdot s(y) = (x \cdot y) + y$.
% \end{description}

% Uz mnogo truda mo\v ze se sada (sli\v cno dokazu koga smo videli) pokazati da su sabiranje i mno\v zenje komutativne
% i asocijativne operacije, i da imamo distributivnost mno\v zenja prema sabiranju. Dakle, sve \v sto o\v cekujemo od
% ove dve operacije mo\v ze se pokazati strogo formalno.

% \section{Gedelove teoreme o nepotpunosti}

% Gedelove (G\"odel) teoreme o nepotpunosti (ima ih dve)
% spadaju me\dj u najva\v znije rezultate savremene logike i imaju duboke posledice ne samo u matematici.
% Recimo, Drugi Penrouzov argument o nemogu\'cnosti kognitivne simulacije iz 1994.~se oslanja upravo da Drugu Gedelovu
% teoremu o nepotpunosti.

% Gedelove teoreme govore o dokazivosti u formalnim aksiomatskim teorijama \v ciji aksiomatski sistemi imaju \enquote{pravilnu
% strukturu} koja se mo\v ze ma\v sinski odlu\v citi. Prva Gedelova teorema ka\v ze da u svakom konzistentnom
% formalnom sistemu \v cije aksiome se mogu algoritamski nabrojati i u kome se mo\v ze izvesti ozvesna
% količina aritmetike, postoje iskazi na jezika tog sistema koji se ne mogu ni dokazati ni opovrgnuti u unutar sistema.
% Druga Gedelova teorema onda pokazuje da takav formalni sistem ne mo\v ze dokazati za sebe da je konzistentan
% (pod pretpostavkom da zaista i jeste konzistentan).

% Za formalni sistem $F$ ka\v zemo da je \emph{konzistentan} ako
% $$
  % \text{\textbf{\textit{nije} }} \calF \vdash \bot.
% $$
% Po\v sto znamo za \v cuvenu tautologiju \emph{ex falso quodlibet}:
% $$
  % \bot \Rightarrow \phi, \text{ za svaku formulu } \phi,
% $$
% mo\v zemo, ekvivalentno, re\'ci da je formalni sistem $F$ konzistentan ako
% $$
  % \text{postoji formula $\phi$ takva da \textbf{\textit{nije} }} \calF \vdash \phi.
% $$
% Formalni sistem je \emph{potpun} ili \emph{kompletan} ako za svaku formulu $\phi$:
% $$
  % \text{ili } \calF \vdash \phi \text{ ili } \calF \vdash \lnot\phi.
% $$
% To zna\v ci da u potpunim formalnim sistemima za svaku formulu $\phi$ mo\v zemo utvrdimo njenu ta\v cnost
% \emph{sintaksno}: ili je $\phi$ dokazivo iz $\calF$ (\v sto zna\v ci da je $\phi$ ta\v cno u sistemu $\calF$),
% ili je $\lnot\phi$ dokazivo iz $\calF$ (\v sto zna\v ci da je $\phi$ neta\v cno u sistemu $\calF$).

% Jedna od osnovnih premisa Hilbertovog programa je da je matematika konzistentna i da \'cemo
% sve u matematici mo\'ci da svedemo na formalnu dedukciju.

% Avaj, nije tako! Ranih 1930-ih Gedel je poslao Hilbertov program u arhive istorije matematike svojim
% teoremama o nepotpunosti.

% \begin{THM}[Prva Gedelova teorema o nepotpunosti]
  % Posmatrajmo formalni sistem $F$ koji sadr\v zi sve logi\v cke aksiome, aksiome Peanove aritmetike, kao i
  % formule kojima je definisano sabiranje i mno\v zenje. Takav formalni sistem mo\v ze da sadr\v zi i druge
  % ne-logi\v cke aksiome sve dok postoji algoritam koji mo\v ze sve da ih izlista.
  % Ako je $\calF$ konzistentan, onda nije potpun, \v sto zna\v ci da postoji re\v cenica $\phi$ na jeziku tog sistema
  % takva da ni $\phi$ ni $\lnot\phi$ nisu dokazivi unutar sistema.
% \end{THM}

% Klju\v cna dosetka u Gedelovim dokazima je upotreba aritmetike da bi se kodirale
% formule i dokazi formalne logike. Ako je $F$ formalni sistem kao u formulaciji Prve Gedelove teoreme,
% onda se (uz mnogo muke!) na jeziku sistema $F$ mo\v ze formulisati re\v cenica $\Consist(F)$ koja kodira stav:
% $$
   % \Consist(F) \equiv \text{\enquote{Sistem $F$ je konzistentan}}.
% $$

% \begin{THM}[Druga Gedelova teorema o nepotpunosti]
  % Posmatrajmo formalni sistem $F$ kao u Prvoj Gedelovoj teoremi. 
  % Ako je $\calF$ konzistentan, onda
  % $$
    % \text{\textbf{\textit{nije} }} \calF \vdash \Consist(F).
  % $$
% \end{THM}

% Dakle, ako je formalni sistem dovoljno mo\'can da se u njemu mogu obavljati logi\v cke i aritmeti\v cke operacije,
% ali i dovoljno jednostavan da se sve \v sto on inicijalno zna o sebi mo\v ze izlistati jedim algoritmom,
% onda taj sistem nema dovoljno logi\v cke snage da doka\v ze sopstvenu konzistentnost!

\section{Predikatski ra\v cun u analizi programskog k\^oda}

Sada \'cemo na dva slo\v zenija primera pokazati kako se ve\v stine koje smo do sada stekli koriste u formalnoj analizi programskog k\^oda.

\begin{EX}\label{ft05.ex.max-niza}
  Posmatrajmo slede\'ci programski fragment:
{\small
\begin{verbatim}
    int i, p, m, n;
    // ovde ucitamo n
    int[] A = new int[n];
    // ovde ucitamo elemente niza A
    p = 0; m = A[0];
    // (1)
    for(i = 0; i < n; i++) {
      if(A[i] > m) {
        p = i; m = A[i];
      }
    }
    // (2)
\end{verbatim}
}
\noindent
  Dokazati da po okon\v canju ovog k\^oda promenljiva $m$ sadr\v zi najve\'ci element niza~$A$, a promenljiva $p$ njegovu
  poziciju u nizu. Drugim re\v cima, dokazati da u ta\v cki (2) k\^oda va\v zi slede\'ca formula:
  $$
    0 \le p \le n - 1 \;\;\land\;\; m = A[p] \;\;\land\;\; (\forall j)(0 \le j \le n - 1 \Rightarrow A[j] \le m).
  $$
  (Prva dva uslova ka\v zu da je $p$ indeks niza i da se $m$ nalazi na mestu $p$ u nizu, dok tre\'ci uslov govori da je
  $m$ ve\'ci od svakog elementa niza -- drugim re\v cima, $m$ je maksimum niza.)
\end{EX}
\begin{proof}[Re\v senje]
  Ozna\v cimo sa $\phi(k)$ slede\'cu formulu:
  $$
    0 \le p \le n - 1 \;\;\land\;\; m = A[p] \;\;\land\;\; (\forall j)(0 \le j \le k - 1 \Rightarrow A[j] \le m).
  $$
  Lako se pokazuje da u ta\v cki (1) k\^oda va\v zi formula $\phi(0)$ koja glasi:
  $$
    0 \le p \le n - 1 \;\;\land\;\; m = A[p] \;\;\land\;\; (\forall j)(\underbrace{0 \le j \le -1}_\bot \Rightarrow A[j] \le m).
  $$
  Naime, znamo da u ta\v cki (1) va\v zi da je $p = 0$ i $m = A[0]$. S druge strane, uslov $0 \le j \le -1$ u tre\'cem \enquote{delu} formule
  je neta\v can. Zato je implikacija ta\v cna, pa je i cela formula ta\v cna.
  
  \v Zelimo da poka\v zemo da je po okon\v canju \enquote{for}-ciklusa ta\v cna formula $\phi(n)$. Ideja dokaza se sastoji u tome
  da se poka\v ze slede\'ce: ako je pre izvr\v senja tela ciklusa ta\v cno $\phi(i)$, onda je po izvr\v senju tela ciklusa
  ta\v cno $\phi(i + 1)$. Pretpostavimo da je ta\v cno $\phi(i)$:
  $$
    0 \le p \le n - 1 \;\;\land\;\; m = A[p] \;\;\land\;\; (\forall j)(0 \le j \le i - 1 \Rightarrow A[j] \le m).
  $$
  i posmatrajmo telo ciklusa:
{\small
\begin{verbatim}
      if(A[i] > m) {
        p = i; m = A[i];
      }
\end{verbatim}
}
\noindent
  Ako je $A[i] \le m$, onda iz $\phi(i)$ sledi da je ta\v cno i slede\'ce:
  $$
    0 \le p \le n - 1 \;\;\land\;\; m = A[p] \;\;\land\;\; (\forall j)(0 \le j \le i \Rightarrow A[j] \le m),
  $$
  odnosno, ta\v cno je $\phi(i + 1)$. Posmatrajmo sada drugu mogu\'cnost: $A[i] > m$. Tada je
  $$
    0 \le p \le n - 1 \;\;\land\;\; m = A[p] \;\;\land\;\; (\forall j)(0 \le j \le i - 1 \Rightarrow A[j] \le A[i])
  $$
  (jer je $A[j] \le m < A[i]$) pa sledi da je ta\v cno i slede\'ce:
  $$
    0 \le p \le n - 1 \;\;\land\;\; m = A[p] \;\;\land\;\; (\forall j)(0 \le j \le i \Rightarrow A[j] \le A[i]).
  $$
  Sada je jasno da nakon naredbi
{\small
\begin{verbatim}
        p = i; m = A[i];
\end{verbatim}
}
  \noindent
  va\v zi $\phi(i + 1)$:
  $$
    0 \le p \le n - 1 \;\;\land\;\; m = A[p] \;\;\land\;\; (\forall j)(0 \le j \le i \Rightarrow A[j] \le m).
  $$
  Dakle, kada se poslednji put izvr\v si telo \enquote{for}-ciklusa znamo da \'ce va\v ziti $\phi(n)$:
  $$
    0 \le p \le n - 1 \;\;\land\;\; m = A[p] \;\;\land\;\; (\forall j)(0 \le j \le n - 1 \Rightarrow A[j] \le m),
  $$
  \v sto je i trebalo dokazati.
\end{proof}


\begin{EX}
  Sortiranje formiranjem uzastopnih minimuma (tzv.\ \emph{selection sort}) se mo\v ze realizovati na slede\'ci na\v cin:
{\small
\begin{verbatim}
  static void swap(int[] A, int i, int j) {
    int tmp = A[i];
    A[i] = A[j];
    A[j] = tmp;
  }
  
  static int posMin(int[] A, int i, int n) {
    int m = i;
    for(int j = i + 1; j < n; j++) {
      if(A[j] < A[m]) { m = j; }
    }
    return m;
  }
  
  static void Sort(int[] A, int n) {
    for(int i = 0; i < n; i++) {
      int m = posMin(A, i, n);
      swap(A, i, m);
    }
  }
\end{verbatim}
}
  \noindent
  Dokazati da metod \textit{Sort} sortira niz $A$ \v cija du\v zina je $n$. Drugim re\v cima, dokazati da
  po okon\v canju metoda \textit{Sort} za elemente niza $A$ va\v zi slede\'ce:
  $$
    (\forall p)(\forall q)(0 \le p \le q < n \Rightarrow A[p] \le A[q]).
  $$
\end{EX}
\begin{proof}[Dokaz]
  Primetimo, prvo, da je formula koja opisuje \v cinjenicu da je niz sortiran ekvivalentna slede\'coj:
  $$
    (\forall p)(0 \le p < n \Rightarrow (\forall q)(p \le q < n \Rightarrow A[p] \le A[q])).
  $$
  Naime, formula data u formulaciji problema ka\v ze: \enquote{za sve $p$ i $q$ koji su indeksi niza, ako je $p \le q$ onda je $A[p] \le A[q]$}.
  S druge strane, formula koju smo upravo naveli ka\v ze: \enquote{za svako $p$ koje je indeks niza, $A[p]$ je manji ili jednak od svakog elementa
  niza koji se javlja iza njega}. Jasno je da su ova dva zahteva ekvivalentna.

  Prvo \'cemo kratko opisati pona\v sanje metoda \textit{swap} i \textit{posMin}. Pona\v sanje metoda \textit{swap} je
  najjednostavnije opisati ovako: ako je pre izvr\v senja naredbe
{\small
\begin{verbatim}
      swap(A, p, q);
\end{verbatim}
}
  \noindent
  va\v zilo $A[p] = \alpha$ i $A[q] = \beta$, tada \'ce nakon izvr\v senja ove naredbe va\v ziti
  $A[p] = \beta$ i $A[q] = \alpha$.
  Sli\v cno Primeru~\ref{ft05.ex.max-niza} mo\v zemo pokazati da nakon izvr\v senja metoda $\mathit{posMin}(A, i, n)$
  kao rezultat dobijamo broj $m$ za koga va\v zi slede\'ce:
  $$
    (\forall q)(i \le q < n \Rightarrow A[m] \le A[q]).
  $$
  Pre\dj imo sada na analizu metoda \textit{Sort}. Neka je $\phi(i)$ slede\'ca formula:
  $$
    (\forall p)(0 \le p < i \Rightarrow (\forall q)(p \le q < n \Rightarrow A[p] \le A[q])).
  $$
  Lako se vidi da pre nego \v sto metod \textit{Sort} krene sa izvr\v savanjem \enquote{for}-ciklusa
  va\v zi formula $\phi(0)$ koja glasi:
  $$
    (\forall p)(0 \le p < 0 \Rightarrow (\forall q)(p \le q < n \Rightarrow A[p] \le A[q])),
  $$
  jer je $0 \le p < 0$ neta\v cno, pa je implikacija trivijalno ta\v cna.
  \v Zelimo da poka\v zemo da je po okon\v canju \enquote{for}-ciklusa ta\v cna formula $\phi(n)$. Kao i ranije,
  ideja dokaza se sastoji u tome da se poka\v ze slede\'ce: ako je pre izvr\v senja tela ciklusa ta\v cno $\phi(i)$,
  onda je po izvr\v senju tela ciklusa ta\v cno $\phi(i + 1)$. Pretpostavimo da je ta\v cno $\phi(i)$:
  $$
    (\forall p)(0 \le p < i \Rightarrow (\forall q)(p \le q < n \Rightarrow A[p] \le A[q])),
  $$
  i posmatrajmo telo ciklusa. Nakon naredbe
{\small
\begin{verbatim}
      int m = posMin(A, i, n);
\end{verbatim}
}
\noindent
  promenljiva $m$ sadr\v zi vrednost za koju znamo da zadovoljava slede\'ce:
  $$
    (\forall q)(i \le q < n \Rightarrow A[m] \le A[q]),
  $$
  pa pre poziva metoda \textit{Swap} va\v zi:
  $$
    (\forall p)(0 \le p < i \Rightarrow (\forall q)(p \le q < n \Rightarrow A[p] \le A[q]))
    \;\;\land\;\;
    (\forall q)(i \le q < n \Rightarrow A[m] \le A[q]).
  $$
  U metodu \textit{Swap} promenljive $A[m]$ i $A[i]$ razmenjuju vrednosti, pa zaklju\v cujemo
  da nakon poziva metoda \textit{Swap} va\v zi:
  $$
    (\forall p)(0 \le p < i \Rightarrow (\forall q)(p \le q < n \Rightarrow A[p] \le A[q]))
    \;\;\land\;\;
    (\forall q)(i \le q < n \Rightarrow A[i] \le A[q]).
  $$
  Sada uo\v cavamo da je desni konjunkt zapravo konsekvent implikacije prve formule za $p = i$, pa dobijamo:
  $$
    (\forall p)(0 \le p < i + 1 \Rightarrow (\forall q)(p \le q < n \Rightarrow A[p] \le A[q])),
  $$
  \v sto je ta\v cno formula $\phi(i + 1)$.
  Dakle, kada se poslednji put izvr\v si telo \enquote{for}-ciklusa znamo da \'ce va\v ziti $\phi(n)$:
  $$
    (\forall p)(0 \le p < n \Rightarrow (\forall q)(p \le q < n \Rightarrow A[p] \le A[q]))
  $$
  \v sto je i trebalo dokazati.
\end{proof}

\section{Zadaci}

\ZAD{07.zad.1}{
Sastaviti predikatsku re\v cenicu (pomo\' cu kvantifikatora) koja ka\v ze:
\begin{enumerate}[(a)]
\item Postoji paran ceo broj.
\item Svaki prirodan broj deljiv je sa samim sobom.
\item Postoji ceo broj deljiv sa tri.
\item Svaki prirodan broj deljiv je sa $-1$.
\item Postoji prirodan broj koji nije kvadrat prirodnog broja.
\item Ne postoji najve\' ci prirodan broj.
\end{enumerate}
}

\ZAD{07.zad.2}{
Napisati predikatsku re\v cenicu re\v cima, i odgovoriti na pitanje da li je data re\v cenica ta\v cna:
\begin{enumerate}[(a)]
\item $(\forall x\in \mathbb Q)(\exists y\in \mathbb Q)(y^2=x)$;
\item $(\forall x\in \mathbb N)(\exists y\in \mathbb Z)(x>y\wedge x\mid y)$;
\item $(\forall x\in \mathbb Z)(\exists y\in \mathbb Z)(x>y\wedge x\mid y)$;
\item $(\forall x\in \mathbb Q)(\exists y\in \mathbb Q)(x\cdot y=\sqrt 2)$;
\item $(\forall x\in \mathbb Q^+)(\exists y\in \mathbb Q)(x\cdot y=\sqrt 2)$;
\item $(\forall x\in \mathbb Q)(\exists y\in \mathbb Q)(x+y=\sqrt 2)$.
\end{enumerate}
}

\ZAD{07.zad.3}{
Napisati re\v cima slede\' ce predikatske re\v cenice, odgovoriti na pitanje da li je data re\v cenica ta\v cna i obrazlo\v ziti odgovor:
\begin{enumerate}[(a)]
\item $(\forall x\in \mathbb N)(\exists y\in\mathbb N)(y<x)$;
\item $(\forall x\in \mathbb N)(\exists y\in\mathbb N)(y<x+1)$;
\item $(\exists x\in \mathbb R)(\exists y\in\mathbb R)(y<x\wedge x<y^2)$;
\item $(\forall x\in\mathbb Q)(\exists y\in \mathbb Q)(x\cdot y\in\mathbb Z)$;
\item $(\forall x\in\mathbb Q)(\exists y\in \mathbb Z)(x\cdot y\in\mathbb Z)$;
\item $(\forall x\in \mathbb R)(\exists y\in\mathbb R)(0<x\cdot y\wedge x\cdot y<2)$;
\item $(\forall x\in \mathbb Q)(\forall y\in\mathbb Q)(\exists z\in\mathbb Q)(x\leq z\leq y\vee y\leq z\leq x)$;
\item $(\exists x\in \mathbb Z)(\exists y\in\mathbb Z)(x\mid y\vee y\mid x)$;

\item $(\forall x\in\mathbb Q^+)(\exists y\in \mathbb Q^+)(x< y\wedge y<x^2)$;
\item $(\exists x\in \mathbb Z)(\forall y\in\mathbb Z)(y+x\cdot y=0)$;
\item $(\forall x\in \mathbb Q)(\forall y\in\mathbb Q)(\exists z\in\mathbb Q)(x\cdot z=y)$;
\item $(\exists x\in \mathbb N)(\exists y\in\mathbb N)(x\neq y\wedge x\mid y\wedge y\mid x)$;
\item $(\forall x\in \mathbb Q^+)(\exists y\in\mathbb N)(\frac 1y<x)$;
\item $(\exists x\in \mathbb Z)(\exists y\in\mathbb Z)(x\neq y\wedge x\mid y\wedge y\mid x)$;
\item $(\exists x\in\mathbb Q^+)(\exists y\in \mathbb Q^+)(x^2< y\wedge y<x)$;
\item $(\exists x\in \mathbb Z)(\forall y\in\mathbb Z)(-y=x\cdot y)$.
\end{enumerate}
}

\clearpage

\ZAD{07.zad.4}{
Negirati predikatsku formulu:
\begin{enumerate}[(a)]
\item $(\forall x)\big[  R(x)\Rightarrow(\exists y)\big[ \big((\forall z)S(z)\big)\vee \neg P(y)      \big]     \big];$
\item $(\exists x)\big[\neg S(x,x)\Rightarrow(\exists y)\big[S(x,y)\vee S(y,x)\Rightarrow(\forall z)R(z,z)\big]\big];$
\item $(\exists x)\big[ \big[(\forall y)(R(x,y)\vee S(y))\big]\Rightarrow (\exists z)\neg T(z)   \big];$
\item $(\forall x)\big[ \big[(\forall y)(R(x,y)\wedge S(x))\big]\Rightarrow(\forall z) R(z,x)      \big];$
\item $(\exists x)\big[(\exists z)\big[R(x,z)\wedge (\forall y)S(z,y)\big]\wedge (\forall t)\big[R(x,t)\Rightarrow(\forall u)R(x,u)\big]      \big];$
\item $(\exists x)(\exists y)\big[\big[S(x,y)\vee R(x,y)\big]\Rightarrow (\exists z)\neg\big[S(x,z)\wedge S(z,y)\wedge R(x,y)\Rightarrow T(x,y)\big]    \big];$
\item $(\forall x)\big[(\exists y)S(x,y)\Rightarrow(\forall y)\big[R(y)\vee S(x,y)\big] \big];$
\item $(\forall x)(\exists y)(\exists z)\big[ P(x)\Rightarrow\big(  P(y)\wedge P(z)\wedge\neg Q(y,z )      \big)   \big].$
\end{enumerate}
}

%\ZAD{07.zad.5}{
%Prevesti predikatsku formulu u preneksni oblik, tj. \enquote{izvu\' ci} kvantifikatore na po\v cetak:\footnote{Preneksni oblik nije ra\dj en na predavanju. Mo\v zda ni ovi zadaci onda ne bi trebali biti tu.}
%\begin{enumerate}[(a)]
%\item $(\forall x)\big[  R(x)\Rightarrow(\exists y)\big[ \big((\forall z)S(z)\big)\vee \neg P(y)      \big]     \big];$
%\item $(\exists x)\big[\neg S(x,x)\Rightarrow(\exists y)\big[S(x,y)\vee S(y,x)\Rightarrow(\forall z)R(z,z)\big]\big];$
%\item $(\exists x)\big[ \big[(\forall y)(R(x,y)\vee S(y))\big]\Rightarrow (\exists z)\neg T(z)   \big];$
%\item $(\forall x)\big[ \big[(\forall y)(R(x,y)\wedge S(x))\big]\Rightarrow(\forall z) R(z,x)      \big];$
%\item $(\exists x)\big[(\exists z)\big[R(x,z)\wedge (\forall y)S(z,y)\big]\wedge (\forall t)\big[R(x,t)\Rightarrow(\forall u)R(x,u)\big]      \big];$
%\item $(\exists x)(\exists y)\big[\big[S(x,y)\vee R(x,y)\big]\Rightarrow (\exists z)\neg\big[S(x,z)\wedge S(z,y)\wedge R(x,y)\Rightarrow T(x,y)\big]    \big];$
%\item $(\forall x)\big[(\exists y)S(x,y)\Rightarrow(\forall z)\big[R(z)\vee S(x,z)\big] \big].$
%\end{enumerate}
%}


\ZAD{10.zad.1}{
Na\' ci model slede\' ce predikatske formule.
\begin{align*}
&(\forall x)(P(x)\Rightarrow \neg Q(x))\wedge (\exists x)R(x)\\
&\wedge (\forall x)(R(x)\Rightarrow P(x))\wedge(\exists x)(R(x)\wedge\neg Q(x))\}
\end{align*}
}

\ZAD{10.zad.2}{
Dokazati da slede\' ca formula nije valjana:
\[(\forall x)(A(x)\Rightarrow B(x))\Rightarrow \neg\big((\exists x)A(x)\wedge(\exists x)\neg B(x)\big).\]
}

\ZAD{10.zad.3}{
Data je formula:
\[\neg(\forall x)\big(  P(x)\Rightarrow(\forall y)\big(P(y)\Rightarrow   ( Q(x)\Rightarrow\neg Q(y) ) \vee (\forall z) P(z)  \big)\big).\]
\begin{enumerate}[(a)]
\item Na\' ci jedan model formule.
\item Dokazati da formula nije valjana.
\end{enumerate}
}

\ZAD{10.zad.4}{
Na\' ci modele slede\' cih formula:
\begin{enumerate}[(a)]
\item $(\forall x)(\forall y)(R(x,y)\wedge R(y,x)\Rightarrow x=y)$;
\item $(\forall x)(\forall y)(\forall z)\big(   R(x,y,z)\wedge R(y,z,x)\wedge R(z,x,y)\Rightarrow T(x,y,z)    \big)$;
\item $(\forall x) R(x,x) \wedge (\forall x)(\forall y)(\forall z) \big(R(x,y)\wedge R(y,z)\Rightarrow R(x,z)\big)$\\
$\wedge (\forall x)(\exists y)\neg R(x,y)$.
\end{enumerate}
}

\ZAD{10.zad.5}{
Dokazati da slede\' ce formule nisu valjane:
\begin{enumerate}[(a)]
\item $(\forall x)(P(x)\Rightarrow Q(x))\Rightarrow \neg\big(    (\exists x)P(x)\wedge(\exists x)\neg Q(x)    \big)$;
\item $(\exists x)\big(P(x)\wedge\neg Q(x)\big)\wedge  (\exists x)\big(Q(x)\wedge\neg S(x)\big)\Rightarrow(\forall x)\big(S(x)\Rightarrow P(x)\big)$;
\item $(\forall x)(\exists y)(\exists z)\big( P(x)\Rightarrow\big(  P(y)\wedge P(z)\wedge\neg Q(y,z )      \big)   \big)$;
\item $(\exists x)(\exists y)\big(\big(S(x,y)\vee R(x,y)\big)\\ \Rightarrow (\exists z)\neg\big(S(x,z)\wedge S(z,y)\wedge R(x,y)\Rightarrow T(x,y)\big)\big)$;
\item $(\forall x)\big( \big((\forall y)(R(x,y)\wedge S(x))\big)\Rightarrow(\forall z) R(z,x)      \big)$;
\item $(\exists x)\big(\neg S(x,x)\Rightarrow(\exists y)\big(S(x,y)\vee T(y,x)\Rightarrow(\forall z)R(z,z)\big)\big)$.
\end{enumerate}
}

\ZAD{10.zad.6}{
Na\' ci model za svaku od slede\' cih predikatskih formula:
\setlength{\multicolsep}{\topsep}
%\begin{multicols}{2}
\begin{enumerate}[(a)]
\item $(\forall x)(P(x)\Rightarrow\neg Q(x))\wedge (\exists x) S(x)\wedge\\(\forall x)(S(x)\Rightarrow P(x))\wedge(\exists x)(\neg S(x)\wedge\neg Q(x))$;
\item $(\forall x)(P(x)\wedge Q(x)\Rightarrow S(x))\wedge(\forall x)(T(x)\Rightarrow S(x))\\
\wedge(\exists x)(P(x)\wedge\neg Q(x)) \wedge (\exists x)(Q(x)\wedge\neg P(x)) \wedge (\exists x)(S(x)\wedge\neg T(x))$;
\item $(\forall x)(P(x)\wedge Q(x)\Rightarrow (x))\wedge (\forall x)(S(x)\Rightarrow T(x)) \\ \wedge (\exists x)(P(x)\wedge\neg Q(x)) \wedge (\exists x)(Q(x)\wedge\neg P(x) \wedge (\exists x)(S(x)\wedge\neg T(x))$.
\end{enumerate}
}


\endinput
%%%%%%%%%%%%%%%%%%%%%%%%%%%%%%%%%%%%%%%%%%%%%%%%%%%%%%%%%%%%%%%%%%%%%%%%%

Predikatski ra\v cun kao formalna teorija.

Ta\v cnost i dokazivost (formula je valjana ako i samo ako je dokaziva).

Teorema nepotpunosti za aritmetiku!
